\section{La referencia del ESP32 frente a SP 800-90B}
\label{sec:esp32_90b}

O8 pide una comparación numérica con el ESP32, y la mitad puede cerrarse
sin la placa: el lado de la referencia. La caracterización anterior del
generador del ESP32 hecha por el autor capturó $10^6$ bits de salida de
\texttt{esp\_fill\_random} en cuatro condiciones de funcionamiento ---
Wi-Fi y ADC activos, solo Wi-Fi, solo ADC, ninguno --- y les aplicó la
entropía de Shannon y SP 800-22; no aplicó los estimadores de SP 800-90B,
y esa es la laguna que se cierra aquí. Para este trabajo, el mismo
firmware, recompilado sin cambios, se ejecutó desde arranque en frío en dos
chips del mismo modelo (ESP32-D0WD-V3), identificados por sus direcciones
MAC: dos veces en uno y tres veces en el otro. Con la captura de agosto,
cuya placa no se identificó, eso da seis repeticiones de $10^6$ bits por
condición. Otras dos capturas con la radio apagada, de $3{,}52\times10^6$
y $8{,}38\times10^6$ bits, se tomaron con el firmware verificador del
Capítulo~\ref{ch:ch5_verification}. Las
veintiséis se pasaron por la implementación de referencia del \ac{NIST} de
los estimadores de SP 800-90B~\cite{nist_ea_github, turan_2018_sp80090b}
(versión 1.1.8), por la vía IID y por la vía no IID, en el formato
binario que la norma prescribe para fuentes de bits. La
Tabla~\ref{tab:esp32_90b} recoge el resultado.

\begin{table}[htbp]
\centering
\caption{El generador del ESP32 frente a los estimadores de SP 800-90B \meas{}.
Condiciones A--D: seis repeticiones de $10^6$ bits cada una en dos chips y
la placa de agosto, media y rango sobre las repeticiones; E y F: capturas
únicas de $3{,}5\times10^6$ y $8{,}4\times10^6$ bits. El techo es el
valor más alto que el estimador del valor más común (MCV) puede devolver
para una fuente perfectamente uniforme de esa longitud con un 99\,\% de
confianza; la cifra no IID es el mínimo de los diez estimadores de la
Sección~6.3 de la norma, y la última columna nombra el estimador que la
fijó (veces de seis). Los veredictos de la vía IID se discuten en el texto.}
\label{tab:esp32_90b}
\footnotesize
\begin{tabular}{@{}llrrll@{}}
\toprule
Captura & Condición & $H$ (IID) & Techo & $H$ (no IID) & Limitante \\
\midrule
A & Wi-Fi + ADC & 0,9954 & 0,9963 & 0,844 (0,812--0,877) & Compresión 5, Lag 1 \\
B & Wi-Fi       & 0,9951 & 0,9963 & 0,862 (0,824--0,898) & Compresión 5, LRS 1 \\
C & ADC         & 0,9952 & 0,9963 & 0,836 (0,812--0,867) & Compresión 6 \\
D & ninguna     & 0,9957 & 0,9963 & 0,857 (0,832--0,916) & Compresión 5, Colisión 1 \\
E (este trabajo) & ninguna & 0,9977 & 0,9980 & 0,867 & Compresión \\
F (este trabajo) & ninguna & 0,9982 & 0,9987 & 0,911 & Compresión \\
\bottomrule
\end{tabular}
\end{table}

La resolución va primero, como en todo este trabajo. El estimador MCV
devuelve $-\log_2 p_u$ con $p_u = \hat p + 2{,}576\sqrt{\hat p(1-\hat p)/N}$,
de modo que ni siquiera una fuente perfecta puede puntuar por encima de
0,9963 con $10^6$ bits, de 0,9980 con $3{,}5\times10^6$ ni de 0,9987 con
$8{,}4\times10^6$. Cada una de las veintiséis capturas queda a menos de
0,0035 de ese techo, y ninguna implica un sesgo de unos superior a
$2{,}4\sigma$, que es lo que se observa en veintiséis extracciones de una
fuente sin sesgo.

La vía IID necesita una calibración antes de que puedan leerse sus
veredictos. De las cincuenta y dos ejecuciones (veintiséis capturas en los
formatos de bits y de bytes), cuatro fallan una sola prueba, nunca el
mismo fichero en ambos formatos: una prueba de permutación en cada formato
y dos pruebas $\chi^2$ en el formato de bytes. La frecuencia con que datos
buenos fallan estas pruebas a esa misma longitud no es el $10^{-3}$
nominal, así que se midió, ejecutando la misma herramienta sobre 1000
ficheros de 125\,000 bytes y 200 ficheros de $10^6$ bits del generador del
sistema operativo. En el formato de bytes, la batería de permutación falla
en el 1,9\,\% de los ficheros buenos y el par $\chi^2$ en el 2,4\,\%; en el
formato de bits, la batería de permutación falla en el 2,5\,\% y el par
$\chi^2$ en ninguno de los 200. La cifra de $\chi^2$ se rastreó hasta su
causa reimplementando ambas pruebas, lo que reproduce los cincuenta y dos
valores $p$ de la herramienta hasta la sexta cifra decimal, y ejecutándolas
sobre 5000 ficheros buenos. La prueba de independencia del formato de
bytes reparte 62\,500 pares sin solape entre 65\,536 celdas y agrupa celdas
hasta que cada grupo espera cinco cuentas, pero a esta longitud su $p$ no
es uniforme: la mediana es 0,175 en lugar de 0,5, y cae por debajo de
$10^{-3}$ en el 1,7\,\% de los ficheros, diecisiete veces la tasa nominal.
Las otras tres pruebas $\chi^2$ están bien calibradas, y el defecto
desaparece a partir de $10^6$ bytes. Frente a estas tasas, los fallos del
ESP32 no tienen nada de particular: un fallo de permutación en veintiséis
ejecuciones en formato de bytes donde se esperan 0,49 ($P = 0{,}39$), dos
fallos $\chi^2$ donde se esperan 0,62 ($P = 0{,}13$), y un fallo de
permutación en el formato de bits donde se esperan 0,65 ($P = 0{,}48$). En
un chip, la prueba de independencia de bytes dio $p < 0{,}05$ en seis
capturas de doce, frente a 2,9 esperadas; agrupar las capturas de cada
chip, donde la prueba está calibrada y es más potente, da valores $p$
empíricos de 0,82 y 0,96 en ese chip para pares de bytes alineados con la
palabra de 32 bits y desplazados un byte, y de 0,38 y 0,86 en el otro, así
que no hay dependencia entre bytes vecinos. El suceso más raro es el fallo
de bondad de ajuste de la primera captura de la condición~D
($p = 2{,}4\times10^{-4}$): esa prueba falla en el 0,04\,\% de los ficheros
buenos, así que un fallo así en veintiséis tiene una probabilidad cercana
al 1\,\%. No se reprodujo en las cinco capturas posteriores de la misma
condición ($p = 0{,}32$ a $0{,}96$) ni en E y F, y se registra como azar;
la advertencia de Espressif de que la aleatoriedad del generador no está
garantizada sin la radio o el ADC sigue en pie como una advertencia que
estas pruebas no ven. Formalmente, un único fallo retira la hipótesis IID
para esa captura; como la cifra que se retiene es la no IID, nada de lo
que sigue depende de ello. Leída con su calibración, la vía IID encuentra
la salida del ESP32 indistinguible de una fuente binaria de entropía
completa con este tamaño de muestra. Es el mismo veredicto que dio
SP 800-22, alcanzado por otro camino, y dice lo mismo: no se ve nada mal.

La vía no IID dice otra cosa, y leerla bien importa más que el número. En
veinticinco de las veintiséis capturas, su mínimo lo fija uno de tres
estimadores --- compresión, colisión o subcadena repetida más larga --- con
0,81 a 0,92 bit por bit; en la restante, el predictor Lag, con 0,88; los
estimadores MCV y de Markov nunca bajan de 0,99. Se sabe que esos tres
estimadores son conservadores con datos casi uniformes~\cite{zhu_2017_analysis};
los predictores de Kelsey, McKay y S\"onmez Turan~\cite{kelsey_2015_predictive}
se introdujeron precisamente porque los estimadores anteriores de la norma
juzgan mal fuentes de ese tipo. Las repeticiones muestran que esto es un
comportamiento de los estimadores y no una propiedad de las condiciones.
La estimación por compresión tiene una desviación típica de 0,02 a 0,04
\emph{dentro} de una condición, tan grande como el 0,03 que separa las
medias de las cuatro condiciones, y un análisis de la varianza de un factor
(ANOVA) sobre cada uno de los once estimadores no encuentra diferencia
entre condiciones: el $p$ más pequeño es 0,083, para el predictor
MultiMCW, por encima del 0,05 sin corregir y muy por encima del 0,0045 que
exigen once pruebas. Lo mismo vale entre chips. Agrupando las condiciones,
que no difieren, ningún estimador separa las ocho capturas de un chip y
las doce del otro al nivel corregido; los dos $p$ más pequeños, 0,038 para
Markov y 0,046 para MCV, corresponden a diferencias de 0,0009 bit por bit
junto al techo de resolución, y dos valores por debajo de 0,05 entre once
pruebas es lo que da el azar. El mínimo no IID es 0,845 en un chip y 0,847
en el otro: con dos dispositivos, la referencia no depende de cuál se mida,
a la resolución de estos estimadores. La captura más grande añade una
segunda muestra del efecto de la resolución: en el formato de bytes, su
estimador limitante es la prueba de $t$-tupla con $t = 2$, que devuelve
7,39 bits por byte porque $10^6$ bytes contienen demasiado pocos de los
65\,536 pares posibles para que el más frecuente sea representativo,
exactamente como le ocurre al techo del MCV con $t = 1$. Ningún estimador
resuelve con seis repeticiones diferencia alguna entre Wi-Fi encendido,
ADC encendido y ninguno de los dos, y no se afirma ninguna.

La prueba de reinicios de la Sección~3.1.4 de la norma responde a lo que
la estadística anterior no puede: si un generador no tiene ruido propio,
su salida tras un reinicio en frío debe repetirse. Un firmware dedicado
leyó 1024~bytes con \texttt{esp\_fill\_random} como primera acción de la
aplicación, con la radio y el ADC apagados (condición~D), y el chip de MAC
a0:dd:6c:0e:fb:4c se reinició 1000 veces por su patilla de habilitación
(EN); la ROM de arranque notificó un reinicio por encendido todas las
veces. Dos variantes comparten la aplicación y difieren solo en el
bootloader de segunda etapa: el de serie, que inyecta ruido del SAR ADC en
el generador durante unos 126~ms mientras se ejecuta, y el mismo
bootloader con esa llamada desviada a una función vacía en tiempo de
enlazado, lo que confirma el desensamblado. Otras tres variantes hacen una
segunda lectura tras arrancar el Wi-Fi como en las condiciones A y B, sin
asociarse a ninguna red, y esperar 200~ms, con cada uno de los dos
bootloaders; la tercera es un control que espera el mismo tiempo sin tocar
la radio, de modo que la lectura cae en el mismo tick del planificador. La
Tabla~\ref{tab:esp32_restart} recoge el resultado.

\begin{table}[htbp]
\centering
\caption{Prueba de reinicios del generador del ESP32 \meas{}: 1000 reinicios
en frío por variante, un chip. $H_r$ y $H_c$ son el mínimo de los
estimadores no IID sobre el conjunto de filas y el conjunto de columnas de
la norma (1000 reinicios $\times$ 1000 bits); la prueba falla cuando
cualquiera de los dos queda por debajo de la mitad de la estimación
inicial $H_I$, la cifra no IID más baja de este chip en la misma condición
(0,838 bit por bit en D, 0,859 en B). El formato de bytes da los mismos
veredictos. Los solapes son reinicios cuya salida aparece desplazada
dentro de la de otro, todos idénticos palabra por palabra. Las lecturas
tardías caen entre 352 y 412~ms después de que arranque la aplicación, la
mayoría a 352~ms.}
\label{tab:esp32_restart}
\footnotesize
\begin{tabular}{@{}llrrrl@{}}
\toprule
Bootloader & Lectura & $H_r$ / $H_c$ & Distintas & Solapes & Veredicto \\
\midrule
De serie & en el arranque, radio apagada & 0,898 / 0,854 & 1000 & 0 & pasa \\
Sin ruido de arranque & en el arranque, radio apagada & 0,024 / 0,813 & 759 & 5424 & falla \\
Sin ruido de arranque & tardía, radio apagada (control) & 0,024 / 0,808 & 760 & 5632 & falla \\
Sin ruido de arranque & tardía, Wi-Fi encendido & 0,831 / 0,845 & 1000 & 0 & pasa \\
De serie & tardía, Wi-Fi encendido & 0,844 / 0,807 & 1000 & 0 & pasa \\
\bottomrule
\end{tabular}
\end{table}

Con el bootloader de serie el generador pasa en ambos formatos, y las
1000 salidas difieren entre sí exactamente como lo harían datos
aleatorios: la distancia de Hamming media entre dos reinicios es de 4096,0
bits de 8192, frente a $4096 \pm 45$ esperados. Sin el ruido de arranque
falla de forma tajante, y los datos muestran por qué. Cada reinicio es una
ventana sobre una única secuencia fija, leída desde un punto de partida
ligeramente distinto: dondequiera que el comienzo de la salida de un
reinicio aparece dentro de la de otro, desplazado unas pocas palabras de
32 bits, todo el solape coincide palabra por palabra, 5424 veces sin una
sola diferencia. Lo único que cambia entre reinicios es el instante en que
la aplicación empieza a leer, que varía en unos pocos miles de ciclos de
reloj de un arranque al siguiente. Eso vale entre 6,5 bits por reinicio,
la cota de la norma a partir de la salida más repetida, y 10,8 bits si
todos los puntos de partida fueran igual de probables, de 8192 bits leídos.
Con la radio y el ADC apagados el generador no tiene ruido propio, y la
afirmación de Espressif de que entonces su salida es solo
pseudoaleatoria~\cite{espressif_idf_random} queda confirmada directamente,
cosa que ninguna prueba estadística sobre un flujo continuo podría hacer.
De ello se siguen dos consecuencias. Las capturas de la condición~D de la
Tabla~\ref{tab:esp32_90b} son la salida de un generador determinista
sembrado durante el arranque: los estimadores les acreditan 0,85 bit por
bit, tanto como con la radio encendida, mientras que su entropía física es
como mucho la que inyectó el bootloader, que esos estimadores no miden. Y
una aplicación que deriva una clave justo después del arranque con la
radio apagada depende por completo del ruido del bootloader de serie; sin
él, la clave podría adivinarse en unos pocos miles de intentos, el
problema del arranque en frío descrito por Blanco-Romero
et al.~\cite{blancoromero_2026_entropy}.

Las filas con radio responden a la otra mitad de la afirmación de
Espressif. Sin el ruido de arranque, arrancar el Wi-Fi basta para que la
prueba pase: las 1000 salidas son todas distintas, no se comparte ninguna
ventana, y los 289 pares de reinicios que leían valores idénticos antes de
arrancar la radio leen 289 valores diferentes después. La temporización
por sí sola no lo explica. El control espera el mismo tiempo y lee en el
mismo tick del planificador, con un instante de lectura que varía otro
tanto (3824 frente a 5260 ciclos de reloj entre los reinicios que
comparten el tick), y se repite exactamente como en el arranque: 760
salidas distintas, 5632 solapes idénticos, unos 1700 estados. Con tan
pocos estados, los 797 reinicios que comparten un tick en la variante con
radio mostrarían unas 180 repeticiones solo por temporización; no muestran
ninguna. Sin la radio, el generador no recibe ruido en sus primeros
350~ms de funcionamiento; con la radio, sí. Lo que la prueba no da es
cuánta entropía por bit aporta la radio: los 0,83 bit por bit validados
son la cifra de los estimadores, con su sesgo habitual, no una medida del
ruido de la radio. El autor no ha encontrado esta prueba publicada antes
para el ESP32.

De ello se siguen tres consecuencias para O8. Primera, la cifra que este
TFM retiene como referencia del ESP32 es la no IID, de 0,81 a 0,92 bit por
bit con una media de 0,85 sobre las veinticuatro capturas de las
condiciones A--D, porque es lo que un evaluador acreditaría sin conceder
la hipótesis IID, con la prueba de reinicios como recordatorio de que para
la condición~D es una cifra estadística y no entropía física; los bits en
bruto del motor, cuando la campaña de la
Sección~\ref{sec:protocol} los produzca, deben pasarse por la misma
batería con el mismo $N$, y compararse estimador por estimador, ya que el
mismo sesgo conservador actuará sobre ellos; la calibración de las pruebas
IID a esa longitud se les aplica sin cambios. Segunda, la cifra de Shannon
de la caracterización anterior, 0,99999 bit por bit, y las cifras de
min-entropía de más arriba no se contradicen --- $H_\infty \le H$
siempre --- pero son monedas distintas: acreditar 256 bits de entropía
cuesta 256 bits de salida en la primera y hasta 316 en la segunda, y solo
la segunda es la que puede usar la generación de claves. Tercera, e
independiente de cualquier número, la medida es de una caja negra: la
fuente y el posprocesado que hay tras \texttt{esp\_fill\_random} no están
documentados, y los estimadores de SP 800-90B están diseñados para ruido
en bruto, no para una salida que puede estar ya acondicionada; la prueba
de reinicios muestra hasta dónde puede llegar eso, un flujo de salida
pseudoaleatoria que todos los estimadores puntúan como ruido. Lo que el
ESP32 no puede aportar, devuelvan lo que devuelvan los estimadores, es la
cota que la Sección~\ref{sec:opoint} deriva para el motor a partir de una
medida física. Esa asimetría es la comparación.

\section{Balance de incertidumbre}
\label{sec:uncertainty}

Cinco magnitudes sostienen las conclusiones, y cada una tiene un tipo de
incertidumbre distinto y declarado.

\begin{description}
\item[Jitter simulado.] Suelo de resolución 1,38~fs \simu{}, unas 870 veces
por debajo de la menor magnitud medida. La comprobación de escalado frente
a la ley de raíz cuadrada esperada concuerda dentro del 0,5\,\% en el
régimen válido, y el límite del régimen está documentado en la
Sección~\ref{sec:spice_results}. En el \ac{PDK} abierto el suelo es de
0,70~fs, 259 veces por debajo del jitter, y la densidad de ruido ya no se
supone, sino que se toma de los modelos de dispositivo. Lo que sigue siendo
aproximado es cómo se aplica ese ruido: la inyección estacionaria puede
sobrestimar la varianza térmica hasta en un factor dos frente a un modelo
cicloestacionario, y el cableado previo al layout es una estimación de
1~fF. La cifra es, por tanto, un valor simulado con una banda de un factor
dos, no el jitter de un dispositivo fabricado.
\item[$Q$ recuperado.] El estimador reproduce un $Q$ conocido con un sesgo
inferior al 2\,\% y una dispersión del 4,0\,\% sobre ocho semillas con
20\,000 bloques \simu{}. Sin las correcciones aplicadas aquí, un simple
ajuste por recta sesga el resultado en un $-18$\,\%, así que la corrección
no es cosmética.
\item[Área.] Una cota inferior \simu{}: solo celdas, con el diseño aplanado,
sin rutado, relleno ni árbol de reloj. Los factores de crecimiento
publicados de área de celdas a área de núcleo en esta biblioteca son
$\approx 2{,}35$ para un \ac{AES}, y los ahorros de optimización del
Capítulo~\ref{ch:ch6_asic} son \est{} y explícitamente no se presentan
como medidos.
\item[Min-entropía de la referencia del ESP32.] Medida \meas{} en
veintiséis capturas de $10^6$ a $8{,}4\times10^6$ bits, seis por condición,
de dos chips del mismo modelo identificados por dirección MAC y una placa
del mismo modelo que no se identificó. La
cifra de la vía IID está en el techo de resolución del estimador MCV para
ese $N$ (0,9963 con $10^6$ bits), así que es un suelo situado en el techo
más que una medida de la fuente, y los veredictos IID se leen frente a la
tasa de falsos fallos medida de cada prueba a esa longitud, no frente a la
nominal. La cifra no IID la fijan estimadores cuyo sesgo conservador con
datos casi uniformes supera la dispersión entre condiciones; un análisis
de la varianza no encuentra efecto de la condición ($p \ge 0{,}083$ en
todos los estimadores) ni efecto del chip ($p \ge 0{,}038$, frente a
0,0045 tras la corrección). Dos dispositivos muestran que la referencia no
depende de uno de ellos; no acotan la dispersión en una población de
producción. La prueba de reinicios se ejecutó en un chip, en la
condición~D y con el Wi-Fi arrancado pero no asociado; no se probaron la
condición con ADC, el Bluetooth ni una radio asociada. Acota la entropía que inyecta el
bootloader de serie solo por abajo: que no haya ninguna repetición en 1000
reinicios significa más de unos $1{,}7\times10^5$ estados de partida igual
de probables, unos 17 bits, con un 95\,\% de confianza. El tamaño del
estado interno del generador no está documentado.
\item[Min-entropía del dispositivo construido.] Aún no acotada. El modelo,
el estimador y los umbrales de corte están listos; el número requiere la
campaña de la Sección~\ref{sec:protocol}. Aquí no se ofrece ningún valor,
y no debe inferirse ninguno del jitter simulado.
\end{description}

\section{La campaña pendiente}
\label{sec:protocol}

El único resultado que debe aportar el silicio es la min-entropía de los
anillos tal como se construyan, y el procedimiento se fija de antemano
para que no pueda ajustarse después de ver los datos.

\begin{enumerate}
\item \textbf{Confirmar que los anillos oscilan.} Leer la frecuencia de
cada anillo con el contador integrado antes de fiarse de cualquier medida
de entropía. Un anillo eliminado por el sintetizador, o dos anillos que se
han enganchado entre sí, serían de otro modo indistinguibles de una fuente
que funciona.
\item \textbf{Medir el jitter} con $K_D = 1$ por el método de Fischer
y Lubicz~\cite{fischer_2014_embedded}, barriendo la distancia entre pares
$M$ y ajustando la región lineal, con al menos tres anillos para que la
componente global se separe de la local~\cite{lubicz_2024_jitter}.
Ajustar $\sigma^2(t) = at + bt^2$ y quedarse solo con el término térmico.
\item \textbf{Fijar $K_D$} a partir del jitter medido para
$Q \ge 0{,}2286$ con margen, y recalcular los umbrales de corte de los
tests de salud para la $H_\infty$ medida en lugar del 0,98 supuesto, con
\texttt{cortes\_salud.py}.
\item \textbf{Recoger bits en bruto} en ese punto de funcionamiento y
ejecutar los estimadores de \ac{NIST} SP 800-90B, contrastando la
estimación empírica con la cota del modelo. El acuerdo respalda el modelo;
el desacuerdo es un resultado sobre las hipótesis del modelo, y en ambos
casos se informa de ello.
\item \textbf{Repetir con distintas ubicaciones} --- anillos en regiones
separadas frente a anillos adyacentes --- para cuantificar el acoplamiento
sobre el que advierte la bibliografía, que es la medida que la
arquitectura de este motor se eligió para hacer posible.
\end{enumerate}

Hasta que se cierre el paso 4, la afirmación del motor es precisa y
limitada: la cadena desde el ruido físico hasta el material de clave está
diseñada, implementada y verificada frente a las normas, y la cota que
producirá es calculable a partir de una medida que aún no se ha tomado.
